% !TEX root = ../../tir_monograph_v12.tex
% V12_CANONICAL_KAPPA_OWNER
% Canonical theorem source:
% TIR/foundations/TIR_KAPPA_FLAVOUR_MIXING_NORMALIZATION_V0_1.md
% Validator:
% TIR/validation/tir_kappa_flavour_mixing_normalization_v0_1.py
\chapter[The Kappa Normalization]{The \texorpdfstring{$\kappa=\ln2/(24\pi)$}{kappa = ln2/(24 pi)} Normalization}
\label{ch:v12_kappa_normalization}

\paragraph{Established context.}
The binary-information identity $H_2(1/2)=\ln2$ is standard Shannon
information theory \cite{shannon1948,CoverThomas2006}.  The angular closure of
one full turn is $2\pi$, and $\dim\mathfrak{su}(3)=8$ is the standard
special-unitary Lie-algebra dimension.  TIR combines these established
ingredients with the three-flavour carrier derived in the preceding chapter.

\vTwelveStatus{B}{--}{PASS}

\section{Binary information numerator}
For a binary distribution $(1-p,p)$,
\[
H_2(p)=-p\ln p-(1-p)\ln(1-p).
\]
At the balanced relational coordinate,
\begin{equation}
\boxed{
\sigma_\star=\frac12,
\qquad
I_\star=H_2(1/2)=\ln2.
}
\label{eq:v12-kappa-binary-information}
\end{equation}
This supplies the informational numerator.

\section{Mixing-channel multiplicity}
From Chapter~\ref{ch:v12_three_flavour_su3},
\[
N_F=3,
\qquad
\dim_{\mathbb R}\mathfrak{su}(3)_F=8.
\]
The generator--flavour incidence set therefore has
\begin{equation}
\boxed{
N_{\rm mix}
=N_F\dim\mathfrak{su}(3)_F
=3\cdot8
=24.
}
\label{eq:v12-kappa-mixing-count}
\end{equation}
Equivalently,
\[
N_{\rm mix}
=N_F(N_F^2-1)
=3(3^2-1)
=24.
\]

\section{Half-turn phase unit}
The primitive half coordinate acts on the standard angular closure as
\begin{equation}
\boxed{
\Delta\phi_{1/2}
=\frac12(2\pi)
=\pi.
}
\label{eq:v12-kappa-half-turn}
\end{equation}
Assigning one such phase unit to each primitive generator--flavour incidence
channel gives the total mixing-phase measure
\begin{equation}
\boxed{
\Phi_{\rm mix}
=N_{\rm mix}\Delta\phi_{1/2}
=24\pi.
}
\label{eq:v12-kappa-mixing-phase-measure}
\end{equation}
Thus
\begin{equation}
\boxed{
24\pi
=\underbrace{(3^2-1)}_{8\ \mathrm{mixing\ directions}}
\underbrace{3}_{\mathrm{flavours}}
\underbrace{\pi}_{\mathrm{half\!\!-turn\ phase}}.
}
\label{eq:v12-kappa-denominator-factorization}
\end{equation}

\section{Structural normalization}
The TIR information-per-mixing-phase normalization is
\begin{equation}
\boxed{
\kappa
=\frac{I_\star}{\Phi_{\rm mix}}
=\frac{\ln2}{N_F(N_F^2-1)\pi}
=\frac{\ln2}{3(3^2-1)\pi}
=\frac{\ln2}{24\pi}.
}
\label{eq:v12-kappa-canonical}
\end{equation}
Numerically,
\[
\kappa
=0.009193150006360484\ldots .
\]
Equation~\eqref{eq:v12-kappa-canonical} is the canonical v12 publication
location of the full derivation.  Downstream chapters refer to this equation and
to its theorem/validator provenance.

\section{Independent tetrahedral integer crosscheck}
The spatial branch obtains the regular tetrahedral finite cell with abstract
automorphism group
\[
\operatorname{Aut}(\Delta^3)\cong S_4,
\qquad
|S_4|=24.
\]
Hence
\begin{equation}
\boxed{
3\,\dim\mathfrak{su}(3)_F
=3\cdot8
=24
=|S_4|.
}
\label{eq:v12-kappa-tetra-crosscheck}
\end{equation}
The flavour-mixing incidence route supplies the normalization dependency; the
spatial route supplies an independent finite-symmetry occurrence of the same
integer.

\section{Exact phase-rate consequence}
Let
\[
\omega\equiv\frac{d\phi}{dt}=2\pi f,
\qquad
d\mathcal I\equiv\kappa\,d\phi.
\]
Then
\begin{equation}
\boxed{
\Gamma_{\mathcal I}
\equiv\frac{d\mathcal I}{dt}
=\kappa\omega
=\frac{\ln2}{12}f.
}
\label{eq:v12-kappa-phase-rate}
\end{equation}
One complete angular cycle carries
\begin{equation}
\boxed{
\Delta\mathcal I_{\rm cycle}
=2\pi\kappa
=\frac{\ln2}{12}.
}
\label{eq:v12-kappa-information-cycle}
\end{equation}

\section{Constraint manifold}
For
\[
\mathbf q=(\kappa,\omega,f,\Gamma_{\mathcal I})
\]
and constraints
\[
C_1=\kappa-\frac{\ln2}{24\pi}=0,
\qquad
C_2=\omega-2\pi f=0,
\qquad
C_3=\Gamma_{\mathcal I}-\kappa\omega=0,
\]
the constraint Jacobian has rank three.  The declared subsystem is therefore
one-dimensional and can be parameterized by
\begin{equation}
\boxed{
\mathbf q(f)=
\left(
\frac{\ln2}{24\pi},
2\pi f,
f,
\frac{\ln2}{12}f
\right).
}
\label{eq:v12-kappa-constraint-curve}
\end{equation}

\section{Reproducibility}
The theorem source is
\path{TIR/foundations/TIR_KAPPA_FLAVOUR_MIXING_NORMALIZATION_V0_1.md}.
Its deterministic numerical/algebraic validator is
\path{TIR/validation/tir_kappa_flavour_mixing_normalization_v0_1.py}.
